\SourceLeaf{379}%
\needspace{12\baselineskip}
\begin{center}[3.]\end{center}

Die allgemeinste Auflösung der Aufgabe, alle Gestalten zu bestimmen, die eine gegebene krumme Fläche annehmen kann, beruht auf folgendem. Es sei $P$ der Punkt der Himmelskugel, welchem die Normale auf der Oberfläche der Superficies in einer Gestalt entspricht, $Q$ der demselben Punkt in einer andern Gestalt entsprechende; eben so correspondiren die Punkte $P',P'',P'''$ etc. für eine Gestalt den Punkten $Q',Q'',Q'''$ [etc.] für die andere. Dann besteht die Auflösung darin, dass jede geschlossene Figur $PP'P''P'''\, .\, .\, .$ auf der Kugelfläche mit der correspondirenden $QQ'Q''Q'''\, .\, .\, .$ einerlei Inhalt habe.

\textls[120]{Spätere Anmerkung.} Dieses schöne Theorem ist zwar vollkommen richtig, aber zur Auflösung nicht hinreichend.

\needspace{8\baselineskip}
\begin{center}BEMERKUNGEN.\end{center}

Die Notizen [2] und [3] finden sich in dem Handbuche »Den astronomischen Wissenschaften gewidmet«, das \textsc{Gauss} im November 1801 begonnen hat; sie stammen vermuthlich ungefähr aus derselben Zeit, wie der unter [1] abgedruckte Brief von \textsc{Gauss} an \textsc{Schumacher}.
\begin{flushright}\textsc{Stäckel}.\end{flushright}
